\diarytitle{0122}{曲面 z=4-x^2-y^2 の下側の体積（極座標変換）}

連続関数 $z = f(x,y) \ge 0$ のグラフと平面 $z=0$、および底面 $D$ とで囲まれる立体の体積は
\[
V = \iint_{D} f(x,y)\,dx\,dy
\]
で与えられる。これは、$D$ を微小面積 $dA$ に分割し、各微小面積の上に高さ $f(x,y)$ の柱を立てて体積 $f(x,y)\,dA$ を足し合わせる（リーマン和の極限をとる）ことから従う。回転対称な曲面（例えば $z = f(x^2+y^2)$）の場合は、極座標変換により積分が簡単になることが多い。

$z \ge 0$ となるのは $x^{2}+y^{2} \le 4$ のときであるから、底面領域は
\[
D = \{(x,y) \mid x^{2}+y^{2} \le 4\}
\]
であり、求める体積は
\[
V = \iint_{D} (4 - x^{2} - y^{2})\,dx\,dy
\]
極座標 $x=r\cos\theta,\ y=r\sin\theta$（$0\le r\le 2,\ 0\le\theta\le 2\pi$）を用いると、$x^2+y^2=r^2$ であるから
\[
V = \int_{0}^{2\pi}\int_{0}^{2} (4 - r^{2})\, r\,dr\,d\theta
= \int_{0}^{2\pi}\int_{0}^{2} (4r - r^{3})\,dr\,d\theta
\]
内側の積分は
\[
\int_{0}^{2} (4r - r^{3})\,dr = \left[2r^{2} - \frac{r^{4}}{4}\right]_{0}^{2} = (8 - 4) - 0 = 4
\]
であるから
\[
V = \int_{0}^{2\pi} 4\,d\theta = 4 \cdot 2\pi = 8\pi
\]
よって
\[
\boxed{\; V = 8\pi \;}
\]
（検算: この立体は放物面 $z=4-r^2$（$0\le r\le2$）を底面とする回転体であり、円柱殻の方法でも
$V = \int_0^2 2\pi r (4-r^2)\,dr = 2\pi\left[2r^2 - \dfrac{r^4}{4}\right]_0^2 = 2\pi\cdot4 = 8\pi$
と一致し、答えが正しいことが確認できる。）
